\section{生成函数把递推变成代数}
\label{sec:04-Discrete-Mathematics/02-Recurrences/03}

给你一条递推公式 \(F_n=F_{n-1}+F_{n-2}\)，初值 \(F_0=0,F_1=1\)，问第 100 项是多少。你当然可以从 \(F_2\) 开始逐项算到 \(F_{100}\)——只是 98 步加法，计算机不嫌累，但你手里只有纸笔的时候，这 98 步每一步都在等上一步的结果完工，中间抄错一个数后面全白干。

更麻烦的是，你想要的不是 \(F_{100}\) 这一个数。你想知道这个序列增长多快，想知道第 \(n\) 项能不能写成一个关于 \(n\) 的封闭表达式，想知道当 \(n\) 很大时相邻项的比值趋于什么。逐项递推给不了这些答案——它只告诉你下一项，不告诉你整条序列的形状。

于是你想到上一节的特征根方法：假设解形如 \(\lambda^n\)，代入递推得到特征方程，解出通项。这条路确实通向 Binet 公式。但特征根法依赖一个隐含假设——解空间是低维线性空间，通解是几个指数函数的线性组合。如果递推不是常系数线性的呢？如果递推涉及组合拆分（``规模 \(n\) 的对象由规模 \(k\) 和 \(n-k\) 的两部分拼成''），指数假设还有用吗？

\subsection{换一个视角：把整条序列装进一个对象}

逐项递推的问题在于你盯着单个下标。特征根法的局限在于你假设了解的形式。有没有办法同时操纵整条序列——对所有 \(n\) 一起做运算？

试试看。把序列 \((a_n)_{n=0}^{\infty}\) 的所有项打包进一个形式幂级数：

\[
A(x)=\sum_{n=0}^{\infty}a_nx^n.
\]

这里的 \(x\) 不代入数值，它只是一个记号——\(x^n\) 的指数标记着这是第 \(n\) 项。说它是``形式''幂级数，意思是目前只关心系数的代数运算，不讨论 \(x\) 取什么值时级数收敛。

这个打包带来的第一个好处：下标平移变成简单的代数操作。乘以 \(x\) 会让每一项的指数加 1，对应系数整体后移一位：

\[
xA(x)=\sum_{n=0}^{\infty}a_nx^{n+1}=\sum_{n=1}^{\infty}a_{n-1}x^n.
\]

再乘一次 \(x\)，系数后移两位。于是递推中的 \(a_{n-1}\)、\(a_{n-2}\) 这些带偏移的下标，在生成函数的世界里只不过是乘了 \(x\) 或 \(x^2\)。

\subsection{用 Fibonacci 演一遍：递推变成方程}

设 \(F_0=0,F_1=1\)，对 \(n\ge 2\) 有 \(F_n=F_{n-1}+F_{n-2}\)。定义

\[
F(x)=\sum_{n=0}^{\infty}F_nx^n.
\]

现在把递推关系在 \(n\ge2\) 上逐项乘以 \(x^n\) 并求和：

\[
\sum_{n\ge2}F_nx^n=\sum_{n\ge2}F_{n-1}x^n+\sum_{n\ge2}F_{n-2}x^n.
\]

左边差前两项就是完整的 \(F(x)\)：\(\sum_{n\ge2}F_nx^n=F(x)-F_0-F_1x=F(x)-x\)（因为 \(F_0=0\)）。

右边第一项提出一个 \(x\)：\(\sum_{n\ge2}F_{n-1}x^n=x\sum_{n\ge2}F_{n-1}x^{n-1}=x\sum_{m\ge1}F_mx^m\)。而 \(\sum_{m\ge1}F_mx^m=F(x)-F_0=F(x)\)（再次用到 \(F_0=0\)）。所以这一项就是 \(xF(x)\)。

右边第二项提出 \(x^2\)，同理得到 \(x^2F(x)\)。

于是递推变成了一个关于 \(F(x)\) 的代数方程：

\[
F(x)-x=xF(x)+x^2F(x).
\]

解出：

\[
F(x)=\frac{x}{1-x-x^2}.
\]

一条原本需要逐项迭代的递推，现在缩成了一个有理函数。所有 \(F_n\) 的信息都编码在这个分式的系数里。

\subsection{从有理函数取回通项}

有了 \(F(x)\) 的封闭形式，下一步是把它展开成幂级数，读出 \(x^n\) 的系数。分母 \(1-x-x^2\) 是二次式，做部分分式分解：

\[
\frac{x}{1-x-x^2}=\frac{1}{\sqrt{5}}\left(\frac{1}{1-\varphi x}-\frac{1}{1-\psi x}\right),
\]

其中 \(\varphi=\frac{1+\sqrt{5}}{2}\)，\(\psi=\frac{1-\sqrt{5}}{2}\) 正是上一节特征方程的两个根。每个分式都是几何级数的生成函数 \(\frac{1}{1-cx}=\sum_{n=0}^{\infty}c^nx^n\)，展开后提取 \(x^n\) 系数即得 Binet 公式：

\[
F_n=\frac{1}{\sqrt{5}}\left(\varphi^n-\psi^n\right).
\]

两条路线在这里汇合了——特征根法从解的指数形式出发，生成函数法从代数操作出发，最终到达同一个表达式。但生成函数法多做了一件事：它在得到通项之前，先让你看到了序列的整体结构——\(F(x)\) 这个有理函数说明 Fibonacci 序列的生成函数是``有理''的，分母的次数等于递推阶数，分子的系数由初值决定。这不是巧合，任何常系数线性递推的生成函数都是有理函数。

\subsection{乘积对应``总规模怎样拆分''}

生成函数的第二个核心操作来自 Cauchy 乘积。若 \(A(x)\) 的系数 \(a_n\) 数第一类规模为 \(n\) 的对象，\(B(x)\) 的系数 \(b_n\) 数第二类，则乘积

\[
A(x)B(x)=\sum_{n=0}^{\infty}\left(\sum_{k=0}^{n}a_k b_{n-k}\right)x^n
\]

的 \(x^n\) 系数恰好枚举了``从第一类取规模 \(k\)、从第二类取规模 \(n-k\)，拼成总规模 \(n\) 的有序对''的所有方案。

这个卷积结构直接对应组合学中的加法原理和乘法原理：总规模 \(n\) 在所有可能的 \(k\) 上拆分为两部分，每部分的方案数相乘，不同 \(k\) 的方案求和。递推中的``\(a_n\) 等于若干个 \(a_{n-k}\) 的线性组合''之所以能翻译成生成函数方程，就是因为下标偏移（乘以 \(x^k\)）与卷积恰好对齐。

举一个具体例子。假设你有 1 元、2 元、5 元三种硬币，每种数量不限，问凑出总额 \(n\) 元的方案数 \(c_n\)。你可以任意使用面值 1 元的硬币——用 0 枚、1 枚、2 枚……对应生成函数因子

\[
1+x+x^2+\cdots=\frac{1}{1-x}.
\]

同理，面值 2 元的因子是 \(\frac{1}{1-x^2}\)，面值 5 元的因子是 \(\frac{1}{1-x^5}\)。三种硬币的选择互相独立，总方案数的生成函数是三个因子的乘积：

\[
C(x)=\frac{1}{(1-x)(1-x^2)(1-x^5)}.
\]

展开这个有理函数，\(x^n\) 的系数就是 \(c_n\)。你不需要枚举所有方案，只需要做多项式乘法或有理函数的部分分式。

如果问题条件改变——比如要求考虑硬币的使用顺序、每种硬币最多用 \(m\) 枚——对应的生成函数因子会随之改变（有序对应不同的乘积结构，上限对应截断的几何级数 \(1+x^k+\cdots+x^{mk}\)）。生成函数不会替你决定用什么因子，它只忠实地执行你从对象模型翻译过来的代数。

\subsection{形式运算与解析运算的边界}

形式幂级数的运算只要求：计算任何一个特定系数时，涉及的加法和乘法只有有限项。这足以支撑所有组合恒等式——你不需要知道级数是否收敛，因为你不代入 \(x\) 的值。

但如果你想知道 \(n\to\infty\) 时 \(a_n\) 的渐近行为，或者想对生成函数求导、积分后代入特定的 \(x\) 值，你就跨过了形式运算的边界，进入了复分析的领域。此时收敛半径不能再忽略：\(\sum a_nx^n\) 只在 \(\abs{x}<R\) 内定义了一个真正的函数，\(R\) 由系数的增长速度决定。

Fibonacci 生成函数的收敛半径是 \(1/\varphi\approx 0.618\)。在收敛圆内，你可以自由地做分析操作；到了圆外，级数发散，有理函数表达式虽然形式上仍有定义，但它不再是幂级数的和。

需要代入数值、使用留数定理或做渐近估计时，确认收敛半径。只做系数恒等式时，保持在形式幂级数的框架内——更少的前提，同样有效的结论。

从递推到生成函数，再到部分分式和通项提取，这一整套流程是组合数学的标准管线。下一篇将进入图论：在那里，``对象和关系''不再被编码为序列的下标偏移，而是显式地画成顶点与边。
